\chapter{Divisibility and primes}\label{ch:div}
\index{divisibility}

Number theory begins with one relation: $a \mid b$, ``$a$ divides $b$''.
This chapter builds the two machines it powers --- the Euclidean algorithm
(the fastest proof in mathematics: it \emph{is} an algorithm and a proof of
correctness at once) and unique factorisation --- and ends with the first
tool an olympiad problem reaches for: solving equations in integers. After
this chapter you will be able to run Euclid by hand in seconds, extract
B\'ezout coefficients from it, prove the digit criteria contestants quote
without thinking, and solve any $ax + by = c$ completely.

% from-zero:begin
\begin{rulebox}[Before you start]
Chapters~2--3 (proofs and induction). \emph{Recap:} $a\mid b$ (``$a$ divides $b$'') means $b=ak$ for some integer $k$. A prime is an integer greater than 1 whose only positive divisors are 1 and itself. \textbf{Book~1, Chapter~14} turns this chapter into fast code.
\end{rulebox}
\begin{wins}[After this chapter you can]
\begin{itemize}
\item divide with remainder and use the Euclidean algorithm to compute a gcd by hand;
\item find integers $x,y$ with $ax+by=\gcd(a,b)$ (B\'ezout);
\item factor into primes and use unique factorisation in a proof;
\item solve a linear Diophantine equation or show it has no solution.
\end{itemize}
\end{wins}
\begin{principle}[How to think about this]
Numbers hide structure: every integer is built from primes, and the gcd is the shared part of two numbers. The Euclidean algorithm works because the common divisors of $a$ and $b$ do not change when you replace $a$ by its remainder; once you see \emph{why} nothing is lost, the algorithm stops being a recipe and becomes obvious.
\end{principle}
% from-zero:end

\section{Divisibility and the division algorithm}\label{sec:div}
\index{a|b@$a \mid b$}

\begin{definition}
For integers $a, b$ with $a \ne 0$, \emph{$a$ divides $b$} (write
$a \mid b$) means $b = ak$ for some integer $k$. Then $a$ is a
\emph{divisor}\index{divisor} (or factor) of $b$, and $b$ is a
\emph{multiple} of $a$.
\end{definition}

Note the type: $a \mid b$ is a \emph{statement} (true or false), while
$a / b$ is a \emph{number}. Writing ``$3 \mid 6$'' asserts; writing
``$3 / 6$'' computes. Mixing them is Common Mistake \#1 and it misfires on
every proof that unfolds the definition.

Basic rules, each a two-line element-chase: if $a \mid b$ and $a \mid c$
then $a \mid (b + c)$ and $a \mid (b - c)$; if $a \mid b$ then $a \mid bk$
for any integer $k$; if $a \mid b$ and $b \mid c$ then $a \mid c$. The
engine behind all of them --- and behind half of olympiad number theory ---
is the division algorithm.

\begin{theorem}[Division algorithm]\index{division algorithm}
For integers $a$ and $b$ with $b > 0$, there exist unique integers $q, r$
with $a = qb + r$ and $0 \le r < b$.
\end{theorem}

\begin{proof}
Existence: let $q = \lfloor a/b \rfloor$ (the largest integer with
$qb \le a$) and $r = a - qb$; then $0 \le r < b$ by the floor's definition.
Uniqueness: if $a = q_1 b + r_1 = q_2 b + r_2$ with both remainders in
$[0, b)$, then $(q_1 - q_2)b = r_2 - r_1$; the right side has absolute value
$< b$ and is a multiple of $b$, so it is $0$, giving $r_1 = r_2$ and
$q_1 = q_2$.
\end{proof}

Every statement ``$b \mid a$'' is the special case $r = 0$. The remainder is
the object; Chapter~6 builds arithmetic on it.

\begin{exercise}[easy]\label{exr:divrules}
Prove from the definition: if $a \mid b$ and $a \mid c$, then
$a \mid (3b - 5c)$.
\end{exercise}

\section{Gcd and the Euclidean algorithm}\label{sec:euclid}
\index{gcd}\index{Euclidean algorithm}

\begin{definition}
The \emph{greatest common divisor} $\gcd(a,b)$ of integers $a, b$ (not both
zero) is the largest integer dividing both. If $\gcd(a,b) = 1$, call $a, b$
\emph{coprime}\index{coprime}.
\end{definition}

\begin{lemma}\label{lem:gcdswap}
For $a = qb + r$, $\gcd(a, b) = \gcd(b, r)$.
\end{lemma}

\begin{proof}
Any $d$ dividing $a$ and $b$ divides $r = a - qb$; any $d$ dividing $b$ and
$r$ divides $a = qb + r$. So the two pairs have \emph{identical} divisor
sets, hence identical greatest elements.
\end{proof}

Lemma~\ref{lem:gcdswap} is the whole machine: replace $(a,b)$ by $(b, a
\bmod b)$ until the second entry is $0$; the survivor is the gcd.
Figure~\ref{fig:euclid-rect} shows why it is called an algorithm \emph{and}
a proof: tiling a $252 \times 198$ rectangle greedily with squares, the last
square that fits is the gcd.

\begin{figure}[ht]\centering
% Figure: the Euclidean algorithm as square-tiling of a 252x198 rectangle.
% Source for Figure~\ref{fig:euclid-rect}. TikZ only; regenerable by compiling main.tex.
\begin{tikzpicture}[
  sq/.style={draw=rulegray},
  fit/.style={fill=accentsoft},
  lab/.style={font=\footnotesize\sffamily,text=slate,align=center},
]
% scale: 252 x 198 at 0.012 units -> 3.02 x 2.38
% one 198x198 square, remainder strip 54x198;
% three 54x54 squares, remainder 54x36;
% one 36x36 square, remainder 18x36 = two 18x18
\draw[sq,fit] (0,0) rectangle (2.38,2.38);           % 198x198
\draw[sq] (2.38,0) rectangle (3.03,2.38);            % strip 54 wide
% inside the strip: three 54x54 (height 0.648 each)
\draw[sq] (2.38,0) rectangle (3.03,0.648);           % 54x54 #1
\draw[sq] (2.38,0.648) rectangle (3.03,1.296);       % #2
\draw[sq,fit] (2.38,1.296) rectangle (3.03,1.944);   % 54x36 = 36x36 + 18x36
\draw[sq] (2.38,1.944) rectangle (3.03,2.38);        % leftover sliver (fits remainder)
% labels
\node[lab] at (1.19,1.19) {$198 \times 198$};
\node[lab] at (2.705,0.324) {\scriptsize $54$};
\node[lab] at (2.705,0.972) {\scriptsize $54$};
\node[lab] at (2.705,1.62) {\scriptsize $36$};
\node[lab] at (1.5,2.7) {one big square, then squares filling the remainder};
\node[lab,align=center] at (1.5,-0.55) {last square used is $18 \times 18$ --- and $18 = \gcd(252,198)$};
\end{tikzpicture}

\caption{Euclid, geometrically. Notice what each square size is: the
\emph{side} of each fitted square is a quotient in the algorithm's trace ---
$198, 54, 36, 18$ --- and the smallest tile that tiles the whole rectangle
is the answer.}
\label{fig:euclid-rect}
\end{figure}

\begin{exercise}[easy]\label{exr:gcdcalc}
Run the Euclidean algorithm on $\gcd(1071, 462)$ by hand, then check with
\texttt{bezout.py} (change the two numbers).
\end{exercise}

\section{B\'ezout's identity --- and why the proof is an algorithm}\label{sec:bezout}
\index{Bezout@B\'ezout's identity}

\begin{theorem}[B\'ezout]
For any integers $a, b$ (not both zero), there exist integers $x, y$ with
$ax + by = \gcd(a,b)$.
\end{theorem}

Two proofs, and the useful one is a procedure: run Euclid forward, then read
the trace backwards --- each remainder was itself a remainder, so the gcd
unrolls into a combination. (The slick proof --- smallest positive
combination divides both --- proves existence but hands you no $x, y$;
contest work wants the coefficients.)

\begin{example}\label{ex:bezout}
Compute $\gcd(252, 198)$ and find integers $x, y$ with
$252x + 198y = \gcd(252,198)$.
\par\smallskip
\textbf{Strategy.} Run the algorithm forward, then read it backwards:
each line $a = qb + r$ rewrites as $r = a - qb$, and substituting upward
keeps every term a combination of $252$ and $198$.
\par\smallskip
\begin{solution}
Forward:
\[
252 = 1\cdot198 + 54,\quad
198 = 3\cdot54 + 36,\quad
54 = 1\cdot36 + 18,\quad
36 = 2\cdot18 + 0,
\]
so $\gcd = 18$. Backward:
$18 = 54 - 1\cdot36 = 54 - (198 - 3\cdot54) = 4\cdot54 - 198 =
4\cdot(252 - 198) - 198 = 4\cdot252 - 5\cdot198$.
Check: $4\cdot252 - 5\cdot198 = 1008 - 990 = 18$. So $x = 4$, $y = -5$.
\end{solution}
\end{example}

\begin{checkpoint}[Predict before you run]
Listing~\ref{lst:bezout} mechanises the back-substitution. Predict its
B\'ezout line before scrolling to the output --- and predict the sign
pattern (why is at least one coefficient negative here? --- a small exercise).
\end{checkpoint}

\begin{lstlisting}[language=Python,caption={\texttt{bezout.py}: Euclid
forward, then read backwards (full script in \texttt{code/bezout.py}).},
label={lst:bezout}]
def euclid(a, b):
    steps = []
    while b:
        q, r = divmod(a, b); steps.append((a, q, b, r)); a, b = b, r
    return a, steps
...
A, B = 1, 0                      # deepest: g = 1*a + 0*b
for a, q, b, r in reversed(steps):
    A, B = B, A - q*B            # unwind r = a - q*b
\end{lstlisting}
Real output:
\begin{lstlisting}[numbers=none]
Euclid(252,198):
  252 = 1*198 + 54
  198 = 3*54 + 36
  54 = 1*36 + 18
  36 = 2*18 + 0
gcd = 18
Bezout: 18 = 4*252 + -5*198  (check 18)

Solve 12x + 18y = 30:
  gcd(12,18) = 6, and 6 | 30 -> solutions exist
  particular: x0 = -5, y0 = 5
  general:    x = -5 + 3t,  y = 5 - 2t
  t=-1: x=-8, y=7, check 12x+18y = 30
  t=0: x=-5, y=5, check 12x+18y = 30
  t=2: x=1, y=1, check 12x+18y = 30
\end{lstlisting}

B\'ezout's corollary is the workhorse: if $a \mid bc$ and $\gcd(a,b) = 1$,
then $a \mid c$. (Multiply $ax + by = 1$ by $c$.) One warning: the lemma
\emph{needs} coprimality --- asserting $\gcd(a,b)=1$ without checking it is
Common Mistake \#3.

\begin{exercise}[medium]\label{exr:bezout}
Find integers $x, y$ with $1071x + 462y = \gcd(1071, 462)$, by
back-substitution (do Exercise~\ref{exr:gcdcalc} first).
\end{exercise}

\section{Primes and the fundamental theorem}\label{sec:fta}
\index{prime}\index{fundamental theorem of arithmetic}

\begin{definition}
An integer $p \ge 2$ is \emph{prime} if its only positive divisors are $1$
and $p$; otherwise it is \emph{composite}. ($1$ is neither --- and the
reason is the next theorem.)
\end{definition}

\begin{theorem}[Fundamental theorem of arithmetic]\index{FTA|see{fundamental theorem of arithmetic}}
Every integer $n \ge 2$ has a unique factorisation
$n = p_1^{e_1} p_2^{e_2} \cdots p_k^{e_k}$ into primes (unique up to order).
\end{theorem}

Existence is strong induction (Chapter~3): $n$ is prime or splits
$n = ab$ with $a, b < n$. Uniqueness uses B\'ezout's corollary: if a prime
$p \mid ab$ then $p \mid a$ or $p \mid b$ (Euclid's lemma) --- apply it to
two factorisations and cancel one prime at a time. The poset of
Figure~\ref{fig:poset} draws the divisor world of one number, $60 =
2^2\cdot3\cdot5$.

\begin{figure}[ht]\centering
% Figure: divisibility poset of 60 --- edges are "x | y, nothing between".
% Source for Figure~\ref{fig:poset}. TikZ only; regenerable by compiling main.tex.
\begin{tikzpicture}[
  nd/.style={circle,draw,minimum size=6.4mm,font=\small},
  ed/.style={draw=rulegray},
  lab/.style={font=\footnotesize\sffamily,text=slate},
]
\node[nd] (n1) at (0,0) {1};
\node[nd] (n2) at (-2.2,1.1) {2};
\node[nd] (n3) at (0,1.1) {3};
\node[nd] (n5) at (2.2,1.1) {5};
\node[nd] (n4) at (-3.0,2.2) {4};
\node[nd] (n6) at (-1.0,2.2) {6};
\node[nd] (n10) at (1.0,2.2) {10};
\node[nd] (n15) at (3.0,2.2) {15};
\node[nd] (n12) at (-1.6,3.3) {12};
\node[nd] (n20) at (0,3.3) {20};
\node[nd] (n30) at (1.6,3.3) {30};
\node[nd] (n60) at (0,4.4) {60};
\draw[ed] (n1)--(n2); \draw[ed] (n1)--(n3); \draw[ed] (n1)--(n5);
\draw[ed] (n2)--(n4); \draw[ed] (n2)--(n6); \draw[ed] (n2)--(n10);
\draw[ed] (n3)--(n6); \draw[ed] (n3)--(n15);
\draw[ed] (n5)--(n10); \draw[ed] (n5)--(n15);
\draw[ed] (n4)--(n12); \draw[ed] (n4)--(n20);
\draw[ed] (n6)--(n12); \draw[ed] (n6)--(n30);
\draw[ed] (n10)--(n20); \draw[ed] (n10)--(n30);
\draw[ed] (n15)--(n30);
\draw[ed] (n12)--(n60); \draw[ed] (n20)--(n60); \draw[ed] (n30)--(n60);
\node[lab] at (4.0,1.1) {prime layer};
\node[lab] at (4.0,2.2) {two factors};
\node[lab] at (4.0,3.3) {three};
\end{tikzpicture}

\caption{The divisors of $60$, edges meaning ``divides with nothing
between''. Notice the layering by factor count: every edge multiplies by a
prime --- the lattice \emph{is} the factorisation $2^2\cdot3\cdot5$ drawn as
a graph.}
\label{fig:poset}
\end{figure}

\begin{exercise}[easy]\label{exr:poset}
Draw the divisibility poset of $36$ and of $30$; why does $30$'s have the
shape of a cube while $36$'s does not?
\end{exercise}

\begin{example}\label{ex:taucount}
Count the positive divisors of $360$ without listing them, then count the
positive divisors of any $n = p_1^{e_1}\cdots p_k^{e_k}$.
\par\smallskip
\textbf{Strategy.} A divisor of $360$ is \emph{constrained freedom}: it
uses the same primes with smaller or equal exponents. Each exponent is an
independent choice --- so multiply the choices.
\par\smallskip
\begin{solution}
$360 = 2^3\cdot 3^2\cdot 5$. A divisor has the form $2^a 3^b 5^c$ with
$0 \le a \le 3$, $0 \le b \le 2$, $0 \le c \le 1$: that is
$4\times3\times2 = 24$ divisors (Figure~\ref{fig:divgrid}). In general the
number of positive divisors is
$\tau(n) = (e_1 + 1)(e_2 + 1)\cdots(e_k + 1)$ --- the ``$+1$'' per factor
is the option \emph{take none of that prime}.
\end{solution}
\end{example}

\begin{figure}[ht]\centering
% Figure: divisors of 360 = 2^3*3^2*5 as an exponent grid (4*3*2 = 24 cells).
% Source for Figure~\ref{fig:divgrid}. TikZ only; regenerable by compiling main.tex.
\begin{tikzpicture}[
  cell/.style={draw,minimum size=6.6mm,font=\tiny},
  lab/.style={font=\footnotesize\sffamily,text=slate,align=center}
]
% rows: exponent of 5 (0,1); inside each row a 4x3 grid of (2^a 3^b)
\foreach \r in {0,1} {
  \foreach \a in {0,1,2,3} {
    \foreach \b in {0,1,2} {
      \node[cell] at (0.72*\a, -0.72*\b - 2.4*\r) {};
    }
  }
}
\node[lab] at (1.0,0.75) {exponent of $2$: $0,1,2,3$};
\node[lab,rotate=90] at (-0.85,-0.8) {exponent of $3$: $0,1,2$};
\node[lab] at (3.3,-0.8) {exponent of $5$: $0$};
\node[lab] at (3.3,-3.2) {exponent of $5$: $1$};
\node[lab,align=center]{Each cell is a divisor $2^a 3^b 5^c$:\\
$(3{+}1)\times(2{+}1)\times(1{+}1) = 24$ divisors of $360$.};
\end{tikzpicture}

\caption{The $24$ divisors of $360 = 2^3\cdot 3^2\cdot 5$ as a
$4\times3\times2$ grid of exponent triples --- the count
$(e_1+1)(e_2+1)(e_3+1)$ is the grid's volume.}
\label{fig:divgrid}
\end{figure}

\section{Euclid's infinitude of primes}\label{sec:infprimes}
\index{primes!infinitude of}

\begin{theorem}
There are infinitely many primes.
\end{theorem}

\begin{example}
Prove it.
\par\smallskip
\textbf{Strategy.} Contradiction --- from any finite list, build a number
\emph{no listed prime divides}: multiply them and add one.
\par\smallskip
\begin{solution}
Suppose $p_1, \dots, p_k$ were all the primes. Consider
$N = p_1 p_2 \cdots p_k + 1$. By Theorem~\ref{thm:primefactor} (Chapter~3)
$N$ has a prime factor $p$. But $N \bmod p_i = 1$ for every $i$, so no
listed prime divides $N$: $p$ is a new prime, contradiction.
\end{solution}
\end{example}

Two warnings students always need: the proof does \emph{not} claim
$p_1\cdots p_k + 1$ is prime ($2\cdot3\cdot5\cdot7\cdot11\cdot13 + 1 =
30031 = 59\cdot509$ is composite) --- only that \emph{some} prime factor
escapes the list. And this is a real proof by contradiction, not a
construction: it gives no recipe for the next prime.

\begin{figure}[ht]\centering
% Figure: Sieve of Eratosthenes to 60 --- composites struck, primes boxed.
% Source for Figure~\ref{fig:sieve}. TikZ only; regenerable by compiling main.tex.
\begin{tikzpicture}[
  cell/.style={minimum width=5.6mm,font=\scriptsize},
  pr/.style={font=\scriptsize\bfseries,text=accent},
  st/.style={font=\scriptsize,text=rulegray},
  lab/.style={font=\footnotesize\sffamily,text=slate},
]
% pass 1: every number grey (struck) except 1
\foreach \n in {1,...,60}
{
  \pgfmathtruncatemacro{\rw}{(\n-1)/10}
  \pgfmathtruncatemacro{\cl}{mod(\n-1,10)}
  \node[cell,st] at (\cl*0.58,-\rw*0.42) {\n};
}
% pass 2: the 17 primes <= 60, boxed in accent
\foreach \p in {2,3,5,7,11,13,17,19,23,29,31,37,41,43,47,53,59}
{
  \pgfmathtruncatemacro{\rw}{(\p-1)/10}
  \pgfmathtruncatemacro{\cl}{mod(\p-1,10)}
  \node[cell,pr,fill=white,inner sep=0.5pt] at (\cl*0.58,-\rw*0.42) {\p};
  \draw[accent,rounded corners=1pt] (\cl*0.58-0.24,-\rw*0.42-0.14) rectangle (\cl*0.58+0.24,-\rw*0.42+0.14);
}
\node[lab,anchor=west] at (6.1,-1.0) {grey: struck by 2, 3, 5 or 7};
\node[lab,anchor=west] at (6.1,-1.4) {boxed: the 17 primes $\le 60$};
\node[lab,anchor=west] at (6.1,-2.0) {sieves stop at $\sqrt{60}$ ---};
\node[lab,anchor=west] at (6.1,-2.4) {every composite has a factor $\le\sqrt{n}$};
\end{tikzpicture}

\caption{The sieve to $60$: composites struck, 17 primes boxed. Notice the
proof inside the picture --- every composite $\le 60$ has a prime factor
$\le\sqrt{60}$, so crossing off multiples of $2,3,5,7$ leaves only primes.
That square-root bound is the same lemma that speeds up primality testing
and powers the primality tests of the olympiad track \cite{zhao}.}
\label{fig:sieve}
\end{figure}

\begin{exercise}[medium]\label{exr:inf2}
Prove there are infinitely many primes $p \equiv 3 \pmod 4$. (Copy Euclid:
form $4p_1\cdots p_k - 1$ and check its factors' residues.)
\end{exercise}

\section{Digit criteria and base representations}\label{sec:digits}
\index{divisibility criteria}\index{base representation}

\begin{theorem}[Base representation]\index{base-$b$ representation}
For any base $b \ge 2$, every integer $n \ge 0$ has a unique expansion
$n = d_k d_{k-1}\dots d_1 d_0 = \sum d_i b^i$ with digits
$0 \le d_i < b$ and $d_k \ne 0$.
\end{theorem}

Figure~\ref{fig:base} shows the same number in two clothes. The digit
criteria are polynomial evaluations: writing
$n = \sum d_i 10^i = P(10)$, any criterion comes from what $10$ does modulo
the divisor.

\begin{example}\label{ex:digitrules}
Prove the tests for $9$ and $11$: $9 \mid n$ iff $9 \mid$ digit-sum of $n$;
$11 \mid n$ iff $11 \mid$ alternating digit-sum of $n$.
\par\smallskip
\textbf{Strategy.} Evaluate the polynomial at $10$: since $10 \equiv 1
\pmod 9$, every power $10^i \equiv 1$, so $P(10) \equiv P(1) = $ the digit
sum. Since $10 \equiv -1 \pmod{11}$, powers alternate $+1, -1, +1, \dots$,
giving the alternating sum. The criterion \emph{is} the congruence.
\par\smallskip
\begin{solution}
Write $n = \sum_{i=0}^{k} d_i 10^i$.
\emph{Mod 9:} $10 \equiv 1 \pmod 9$, so $10^i \equiv 1$ and
$n \equiv \sum d_i \pmod 9$; hence $9 \mid n$ iff $9 \mid \sum d_i$.
\emph{Mod 11:} $10 \equiv -1 \pmod{11}$, so $10^i \equiv (-1)^i$ and
$n \equiv d_0 - d_1 + d_2 - \cdots \pmod{11}$; hence $11 \mid n$ iff $11$
divides that alternating sum.
\end{solution}
\end{example}

\begin{figure}[ht]\centering
% Figure: base-b place value --- 457 in base 10 vs base 7.
% Source for Figure~\ref{fig:base}. TikZ only; regenerable by compiling main.tex.
\begin{tikzpicture}[
  cell/.style={draw=rulegray,minimum width=9mm,minimum height=6.5mm,font=\small},
  pow/.style={font=\scriptsize\sffamily,text=slate},
  lab/.style={font=\footnotesize\sffamily,text=slate},
]
% base 10: 457 = 4*100 + 5*10 + 7*1
\node[lab,anchor=east] at (-0.4,0) {base 10:};
\foreach \x/\d/\p in {0/4/{10^2},1/5/{10^1},2/7/{10^0}}
{
  \node[cell] at (\x*0.92,0) {\d};
  \node[pow,anchor=north] at (\x*0.92,-0.42) {$\p$};
}
\node[pow,anchor=west] at (2.1,0.15) {$= 4{\cdot}100 + 5{\cdot}10 + 7$};
% base 7: 457 = 1*343 + 2*49 + 2*7 + 4  -> digits 1 2 2 4
\node[lab,anchor=east] at (-0.4,-1.5) {base 7:};
\foreach \x/\d/\p in {0/1/{7^3},1/2/{7^2},2/2/{7^1},3/4/{7^0}}
{
  \node[cell] at (\x*0.92,-1.5) {\d};
  \node[pow,anchor=north] at (\x*0.92,-1.92) {$\p$};
}
\node[pow,anchor=west] at (3.0,-1.35) {$= 1{\cdot}343 + 2{\cdot}49 + 2{\cdot}7 + 4$};
\node[lab,anchor=west] at (3.0,-2.0) {each digit $<$ base};
\node[lab] at (1.4,-3.0) {the same number $457$, two different clothes};
\end{tikzpicture}

\caption{Place value in two bases. Notice what changes and what doesn't:
the digits differ, the powers differ, the \emph{number} is the same --- base
is a notation, not a property of $n$.}
\label{fig:base}
\end{figure}

Same recipe, every divisor: $3$ works like $9$; $2, 5, 10$ read the last
digit since $10 \equiv 0$; $4$ and $25$ read the last two since
$100 \equiv 0$. Cut-the-Knot collects the whole family~\cite{cuttheknot}.

\begin{exercise}[easy]\label{exr:digit}
Derive the test for $7$ in the style of Example~\ref{ex:digitrules}: what
is $10 \bmod 7$, and why is the classical ``double the last digit, subtract''
rule equivalent?
\end{exercise}

\section{Linear Diophantine equations}\label{sec:dioph}
\index{Diophantine equation}

\begin{theorem}\label{thm:dioph}
$ax + by = c$ has integer solutions iff $\gcd(a,b) \mid c$. If $(x_0, y_0)$
is one solution, all solutions are
$x = x_0 + \frac{b}{\gcd(a,b)}\,t$,\; $y = y_0 - \frac{a}{\gcd(a,b)}\,t$
for $t \in \Z$.
\end{theorem}

\begin{proof}
Direction $\Rightarrow$: $\gcd$ divides $ax + by$. Direction $\Leftarrow$:
B\'ezout gives $ax' + by' = \gcd$; multiply by $c/\gcd$. For the family:
substituting shows every listed pair works; conversely, any two solutions
differ by a solution of $ax + by = 0$, i.e.\ $\frac{a}{\gcd}x =
-\frac{b}{\gcd}y$ --- coprime factors force
$x = \frac{b}{\gcd}t$.
\end{proof}

\begin{example}\label{ex:dioph}
Find all integer solutions of $12x + 18y = 30$.
\par\smallskip
\textbf{Strategy.} Two gates, in order: does $\gcd(12,18) = 6$ divide $30$?
Yes --- so solutions exist. Get one by B\'ezout (or by halving the equation
to $2x + 3y = 5$ and spotting $x = 1, y = 1$); then the family is periodic
with period $b/\gcd = 3$ in $x$.
\par\smallskip
\begin{solution}
Divide through by $6$: $2x + 3y = 5$. Spot $x_0 = 1, y_0 = 1$
($2 + 3 = 5$). By Theorem~\ref{thm:dioph} all solutions are
$x = 1 + 3t$, $y = 1 - 2t$ for $t \in \Z$. (The script's equivalent
parametrisation $x = -5 + 3t$, $y = 5 - 2t$ is the same line re-indexed.)
Check: $12(1 + 3t) + 18(1 - 2t) = 12 + 18 + 36t - 36t = 30$.
\end{solution}
\end{example}

The structure is worth naming: solutions form a \emph{particular solution
plus a lattice} --- a linear Diophantine equation is a line of integer
points, and B\'ezout is how you land on it (Figure~\ref{fig:divline}).

\begin{figure}[ht]\centering
% Figure: Diophantine solutions as equally spaced lattice points on a line.
% Source for Figure~\ref{fig:divline}. TikZ only; regenerable by compiling main.tex.
\begin{tikzpicture}[
  pt/.style={circle,draw,inner sep=1.6pt,minimum size=3mm},
  sol/.style={circle,draw,fill=accent!30,inner sep=1.6pt,minimum size=3mm},
  lab/.style={font=\footnotesize\sffamily,text=slate,align=center}
]
\draw[->,rulegray] (-0.4,0) -- (7.2,0);
\draw[thick,rulegray] (-0.2,-0.75) -- (6.8,2.1);
\foreach \x in {-0.2,0.6,1.4,2.2,3.0,3.8,4.6,5.4,6.2,6.8}
  \node[pt] at (\x,{\x*0.35-0.68}) {};
\node[sol] (s0) at (2.2,0.09) {};
\node[sol] (s1) at (3.8,0.65) {};
\node[sol] (s2) at (5.4,1.21) {};
\node[lab] at (2.2,-0.45) {$x_0$};
\node[lab] at (3.8,1.15) {$x_0 + \tfrac{b}{d}$};
\node[lab] at (5.4,1.71) {$x_0 + 2\tfrac{b}{d}$};
\node[lab] at (3.2,-1.25) {Once one solution exists, all solutions are
spaced $\tfrac{b}{\gcd}$ apart --- an arithmetic family,\\
never a scatter. (For $ax + by = c$.)};
\end{tikzpicture}

\caption{All solutions of $ax + by = c$ sit on one line, spaced
$b/\gcd(a,b)$ apart in $x$ --- an arithmetic family from one seed
$x_0$, never a scatter of independent hits.}
\label{fig:divline}
\end{figure}

\begin{exercise}[medium]\label{exr:dioph}
Solve completely or prove impossible: (a) $15x + 35y = 100$;
(b) $21x + 14y = 63$; (c) $34x + 55y = 1$ (consecutive Fibonacci numbers ---
predict $\gcd$ first).
\end{exercise}

\section{Chapter summary}
\begin{table}[ht]\centering
\begin{tabularx}{\linewidth}{@{}lYY@{}}
\toprule
\tablehead\thead{Object} & \thead{Machine} & \thead{Watch}\\
\midrule
$a \mid b$ & $b = ak$ unfolding & it is a statement, not a fraction\\
$\gcd(a,b)$ & Euclid: $(a,b) \to (b, a \bmod b)$ & log keeps every remainder\\
$ax + by = \gcd$ & B\'ezout by back-substitution & one coefficient usually goes negative\\
prime factorisation & FTA: unique up to order & existence = strong induction\\
$ax + by = c$ & gcd-gate then periodic family & check $\gcd \mid c$ FIRST\\
digit test for $d$ & evaluate $P(10) \bmod d$ & $10 \equiv 1, -1, 0$ cases\\
infinitude of primes & $p_1\cdots p_k + 1$ & shows a \emph{new} prime, not \emph{the} next\\
\bottomrule
\end{tabularx}
\caption{The chapter in one page: every construction doubles as a proof.
The right column is where hand-written solutions actually break.}
\label{tab:div}
\end{table}

\section{Common mistakes}
\index{common mistakes}
\begin{itemize}
\item \textbf{$a \mid b$ read as a fraction.} $a \mid b$ is a proposition;
$a/b$ is a rational number. The error bites when unfolding:
``$3 \mid 6 = 2$'' is nonsense --- the claim is $6 = 3\cdot2$.
\item \textbf{``Prime'' confused with ``odd''.} $9$ and $15$ are odd and
composite; $2$ is prime and even. The definition is about \emph{divisors},
not parity --- most students meet it as a slip in a proof that ``assumes''
a prime is odd.
\item \textbf{Asserting $\gcd(a,b)=1$ unproved.} Euclid's lemma and the
family formula for Diophantine equations both \emph{require} coprimality.
``WLOG $\gcd(p,q)=1$'' is a licence only after you justify it (lowest
terms) or divide it out explicitly.
\end{itemize}

\section{Additional exercises}
\begin{exercise}[easy]
Prove: if $a \mid b$ and $b \mid a$ for positive integers, then $a = b$.
(Where does ``positive'' enter?)
\end{exercise}
\begin{exercise}[easy]
Compute $\gcd(2016, 784)$ and $\gcd(1001, 77)$ by Euclid; confirm with
\texttt{code/bezout.py}.
\end{exercise}
\begin{exercise}[easy]
Prove the divisibility tests for $4$ and $25$ from
Example~\ref{ex:digitrules}'s method.
\end{exercise}
\begin{exercise}[easy]
How many divisors does $360 = 2^3\cdot3^2\cdot5$ have? Prove the formula
$d(n) = \prod (e_i + 1)$ from the FTA.
\end{exercise}
\begin{exercise}[medium]
Prove $\gcd(a, b) \cdot \mathrm{lcm}(a, b) = ab$ for positive $a, b$.
(Track one prime's exponent at a time.)
\end{exercise}
\begin{exercise}[medium]
Prove: if $\gcd(a, b) = 1$ and $a \mid bc$, then $a \mid c$ --- and find a
counterexample when $\gcd(a,b) > 1$.
\end{exercise}
\begin{exercise}[medium]
Show that among any three consecutive integers $> 2$, at least one is
composite; and that $n! + 2, n! + 3, \dots, n! + n$ gives $n - 1$
consecutive composites for any $n$.
\end{exercise}
\begin{exercise}[medium]
Prove a square number never ends in digit $2, 3, 7$ or $8$. (Reduce
$(10a + b)^2 \bmod 10$.)
\end{exercise}
\begin{exercise}[hard]
Prove Euclid's lemma in both directions: $p$ prime iff
$p \mid ab \Rightarrow p \mid a$ or $p \mid b$ --- the reverse direction is
a clean characterisation of primality.
\end{exercise}
\begin{exercise}[hard]
Chickens cost $3$ coins, goats $7$: a farmer spends exactly $200$ coins.
Prove a solution exists and find all solutions ($3x + 7y = 200$, $x, y \ge
0$). Why is this Postage-with-a-budget?
\end{exercise}
\begin{exercise}[hard]
Extend \texttt{code/bezout.py} to print, for each Euclid step, the running
coefficients $(x_i, y_i)$ with $ax_i + by_i = r_i$ maintained as an
invariant --- the \emph{iterative} extended algorithm. Prove its invariant
by induction and verify on $(252,198)$ and $(1071,462)$.
\end{exercise}
\begin{exercise}[easy]
Compute $\tau(2025)$, $\tau(3600)$, and the number of \emph{odd} divisors
of $3600$ --- the odd-divisor count is the same product taken only over odd
primes.
\end{exercise}
\begin{exercise}[medium]
Prove the converse of the gcd-gate for Diophantine equations is tight: if
$d \mid c$ fails, no integer solution exists; then exhibit an
\emph{integers-only} equation with a solution whose real line has
infinitely many integer points ($x - y = 0$) to see why ``one solution''
and ``finitely many'' are unrelated.
\end{exercise}
\begin{exercise}[medium]
A number $n$ is \emph{abundant} if $\sigma(n) > 2n$ where $\sigma$ sums the
divisors. Prove $945$ is the smallest odd abundant number is plausible by
computing $\sigma(945) = \sigma(3^3\cdot5\cdot7)$ using the multiplicative
formula $\sigma(p^e) = (p^{e+1}-1)/(p-1)$, and verify $\sigma(945) = 1920$.
\end{exercise}
\begin{exercise}[hard]
Prove that for every $n \ge 2$ there is a prime strictly between $n$ and
$n!$: adapt Euclid's construction to $N = n! - 1$ --- where does the
argument break if you use $n! + 1$ instead?
\end{exercise}

\section*{Further reading}
\begin{itemize}
\item R.~Hammack, \emph{Book of Proof}, 3rd ed., Ch.~4--5 ---
\url{https://richardhammack.github.io/BookOfProof/} --- free PDF;
divisibility, gcd and Euclid drilled with full proofs.
\item E.~Lehman, F.~T.~Leighton, A.~R.~Meyer, \emph{Mathematics for
Computer Science}, Ch.~9 ---
\url{https://courses.csail.mit.edu/6.042/spring18/mcs.pdf} --- free PDF;
the number-theory chapter with the gcd-gate and Euclid's lemma treated the
way \S\ref{sec:bezout} and \S\ref{sec:dioph} use them.
\item Y.~Zhao, \emph{Modular Arithmetic} ---
\url{https://yufeizhao.com/olympiad/mod2.pdf} --- the handout
that continues this chapter into orders, Fermat--Euler and Wilson.
\item MIT OpenCourseWare, \emph{6.042J} (Fall 2010) ---
\url{https://ocw.mit.edu/6-042JF10} --- number-theory lectures behind
\S\ref{sec:euclid}--\S\ref{sec:fta}.
\item The Open Logic Project ---
\url{https://openlogicproject.org/} --- free texts; set-theoretic
groundwork for treating integers as a structure.
\item Cut-the-Knot, divisibility articles ---
\url{https://www.cut-the-knot.org/} --- the digit-criteria and
sieve-as-proof perspective behind \S\ref{sec:digits}.
\end{itemize}

\begin{rulebox}[Watch instead]
\begin{itemize}
\item \textbf{MIT OCW 6.042J --- ``Lecture 4: Number Theory I''}, 1:20:24,
\url{https://www.youtube.com/watch?v=NuY7szYSXSw} --- Euclid's
algorithm, gcd and divisibility done at the blackboard: \S\ref{sec:euclid}--%
\S\ref{sec:fta}. (Verified 2026-10-11.)
\item \textbf{Khan Academy --- ``Journey into cryptography'' playlist}
\url{https://www.youtube.com/playlist?list=PLSQl0a2vh4HA50QhFIirlEZRXG4yjcoGM}
--- its modular-arithmetic unit covers \S\ref{sec:digits} and the
congruence view. (Verified 2026-10-11.)
\end{itemize}
\end{rulebox}

\chapternotes
The presentation follows MIT 6.042J's number-theory chapter: division
algorithm, Euclid-as-lemma-machine, B\'ezout by back-substitution, and the
gcd-times-lcm identity~\cite{mcs}; Velleman's chapter on divisibility
supplied the WOP-style proofs and the element-chase drills~\cite{velleman},
and Hammack's proof chapters the $\gcd\cdot\mathrm{lcm}$ and
gcd-gate material~\cite{hammack}. The digit-criteria family and the
sieve-as-proof framing follow Cut-the-Knot's divisibility
articles~\cite{cuttheknot}; Zhao's modular-arithmetic handout ladder
(orders, Fermat--Euler, Wilson) is the recommended continuation into
Chapter~6~\cite{zhao}. The contest problems this chapter feeds ---
including AIME~2025~I~P1, AIME~2025~I~P5, AMC~12~2024~A~P4, AMC~12~2024~A~P9,
AMC~12~2024~A~P11, AIME~2024~II~P7 and the classic remainder/digit items
AIME~2018~I~P2, AIME~2022~I~P2, AIME~2020~I~P3, AMC~12~2024~B~P21 --- are on
the AoPS archive~\cite{aops-amc}.
